\section{Coverage and Evaluation Design}
\label{sec:design}

We count coverage first by article and then by issuer, since a single
canonical copy may supply news about several firms. Panel~A of
Table~\ref{tab:coverage} counts each distinct article ID once per row,
while Panel~B follows issuers across the nine complete calendar years
from 2017 through 2025. For a given year an issuer qualifies with
either a dated listing or an
undated current listing supported by public-float evidence in that year
or the preceding year. Qualifying issuers remain in the denominator
even when no articles could be linked to them.

Of 34,781 eligible issuer-years, 33,024 have at least one usable article.
At an annual minimum of 50 articles, 13,918 issuer-years qualify on
usable links, while 8,641 meet the same minimum when each story must
mention the issuer at least three times. The cost of that restriction
also varies with firm size. In the lowest annual quintile of measured
public float, 329 of 6,533 issuer-years (5.0\%) clear the usable-article
minimum and 105 (1.6\%) clear the focused minimum. In the highest
quintile the corresponding counts are 5,790 of 6,526 (88.7\%) and
4,561 (69.9\%). Appendix Table~\ref{tab:size-year} gives the
intermediate quintiles and the breakdown by year.

\begin{table}[t]
  \centering
  \small
  \begin{tabular}{@{}lr@{}}
    \toprule
    Panel A. Article requirements & Article IDs \\
    \midrule
    Staged & 14,807,287 \\
    Consistently dated & 11,287,476 \\
    Usable canonical copy & 11,260,839 \\
    Usable with an accepted issuer link & 5,220,290 \\
    Usable with a focused issuer link & 2,547,597 \\
    \bottomrule
  \end{tabular}

  \smallskip
  \begin{tabular}{@{}lrrr@{}}
    \toprule
    Panel B. Annual minimum & Raw & Usable & Focused \\
    \midrule
    1 article & 33,109 & 33,024 & 31,745 \\
    10 articles & 26,930 & 26,434 & 22,243 \\
    50 articles & 14,539 & 13,918 & 8,641 \\
    \bottomrule
  \end{tabular}
  \caption{Distinct staged IDs surviving the article-level rules
    (Panel~A) and issuer-years meeting each annual article minimum
    during 2017--2025 (34,781 eligible issuer-years in Panel~B). An
    article can count for several issuers
    and qualifies as focused for each linked issuer it mentions at least
  three times.}
  \label{tab:coverage}
\end{table}

To understand which stories enter an issuer's coverage, we examine
first whether the extractor recognizes a company name and then whether
the linker assigns that name to the right issuer. Of the 712
model-reference company spans in 250 sampled leads, GLiNER recovered 412
and Stanza 421 under exact matching, while allowing overlapping boundaries
brings their company-span F1 scores within half a percentage point
(Appendix Table~\ref{tab:app-spans}).

For the identity comparison, we sample extracted spans independently of
the linker's decisions and ask model judges to identify the issuer on
the article date from the marked text and neutral search results, with
the predictions withheld. On held-out extracted spans,
the 475 accepted links have 96.8\% estimated precision (95\% interval
94.0--99.6\%), and recall among reference-linkable mentions is 70.8\%
(64.6--77.1\%). On the same 886 held-out mentions, weighted issuer-link
F1 is 81.81\% with the full linker and 59.46\% when a normalized name
must identify one dated issuer. Restricting learned aliases and their
frequencies to prior calendar years yields 78.57\%
(Appendix Table~\ref{tab:app-link-baselines}).

Having established which issuer a name identifies, we ask whether the
article concerns that issuer. This comparison uses article-level
judgments and is conditional on GLiNER detecting at least one company
in the article. Among the relevance, position and mention-count rules
compared on 217 development articles, the three-mention rule achieved
the highest subject F1. On 183 separate test articles, it retained 37
of 65 judged subject pairs. Relative to counting all accepted links,
requiring three mentions raised subject precision from 31.9\% to 66.1\%
(54.8--77.3\%) but reduced recall from 92.3\% to 56.9\%
(44.3--69.5\%), since some of the excluded stories concerned the issuer.

Because the production panel counts named mentions, a story can fall
short of three even when it continues to discuss the firm using a
pronoun or a description. When we also count nominal and pronominal
mentions in Stanford CoreNLP chains anchored to a single linked issuer
on the same test articles \citep{manning2014corenlp}, twelve additional
article--issuer pairs qualify, seven of which are judged subjects.
Recall rises from 56.9\% to 67.7\% (paired gain of 10.77 percentage
points, 95\% interval [3.92, 19.61]), while precision falls from
66.1\% to 64.7\%. A separate GPT-6.1 Terra review attributed 28 of
the 42 added references in those pairs to the anchored issuer. Seven
of the twelve pairs still qualify using only those confirmed references,
and six of those seven are judged subjects.

Once a story is assigned to a firm, its tone can still differ from
what the news implies for that firm's investors. On the 358 of 400
marked mentions that the primary judge labeled positive, negative or
neutral, GLiNER's company-specific sentiment achieved 46.79\% weighted
macro-F1, compared with 40.84\% when each company received the
Stanford sentiment of its sentence. The paired 95\% interval for the
5.95-point difference spans zero [$-1.33$, $+13.98$]. A blinded
GPT-6.1 Terra review marked 23 of 40 sampled targets as invalid
company spans. We therefore treat issuer sentiment as exploratory.

Similar questions arise when researchers select articles by form or
event. To assess how those choices affect sample composition, blinded
model judges label 400 sampled units in each of three confirmation
frames covering article form, local mention attributes and event
families in input windows. The weighted
press-release share among usable articles is 42.1\% under the
classifier and 4.2\% under the primary judge, with 14 reference-positive
sampled articles. A filter on predicted transactions would miss many
windows the judge labels as transactional, since recall against those
labels is 23.6\%. Appendix
Table~\ref{tab:app-metadata} reports estimates by class.

The unit on which these decisions operate changes across the three
financial applications. For the equity-market volatility tracker we
divide the number of articles containing the required groups of words
by all eligible articles from that publisher in the same month. For a
bank network we also need to know which institutions appear together,
since each eligible article contributes one observation for every
unordered pair of mapped banks it mentions in the declared week. A
three-bank story thus supplies three pairs, while the weight on each
pair counts the stories in which it appears.

For the overnight exercise, identifying the security also determines
which exchange session to use. We pair timely news from two Indian
publishers with US-listed securities, which trade on Nasdaq in 740
article--security rows of the fixed panel and on the NYSE in 364.
For 270 rows the registry lacks a
listing start date, so eligibility also requires usable pre-2021 news
linked to the issuer. To establish what could have been observed before
the open, we require first publication to follow the preceding exchange
close and precede 09:00 New York time, with capture between publication
and the cutoff, inclusive. If the crawler fetches an overnight story
after 09:00, we observe the body only after the cutoff. We measure
drift from the open to the close because the close-to-open reaction
begins before the article appears. The
rolling probes train on returns realized before each refit, while
the encoder was released after the study period and its pretraining
overlap with the archived stories is unmeasured. Each active long or
short leg requires at least two
securities, so a session with predictions contributes zero to the
portfolio mean when neither leg qualifies,
while a day without predictions falls outside the evaluation.
